% Task 8: the text of the task. Tasks.tex puts all tasks into one PDF; the code shown here is in Task8/.
% the buffer figure: producer S1, consumer S2
\newcommand{\bufferfig}{%
\begin{tikzpicture}
  \node[actor=sIndigo, minimum height=0.9cm, minimum width=2.2cm] (p) at (0,0) {producer};
  \foreach \i in {0,1,2} \node[cellused, minimum size=0.75cm] at (3+\i*0.75,0) {};
  \foreach \i in {3,4} \node[cellfree, minimum size=0.75cm] at (3+\i*0.75,0) {};
  \node[actor=sPink, minimum height=0.9cm, minimum width=2.2cm] (c) at (9,0) {consumer};
  \draw[data] (p.east) -- node[above, font=\small, text=nInk] {push} (2.6,0);
  \draw[data] (6.4,0) -- node[above, font=\small, text=nInk] {pop} (c.west);
  \node[font=\small, text=nInk] at (4.5,-0.75) {the buffer (at most $N$)};
  \node[font=\scriptsize, text=nInk] at (3.75,0.65) {used};
  \node[font=\scriptsize, text=nInk] at (5.62,0.65) {free};
\end{tikzpicture}}
\settitle{Task 8}{The video decoder}

Task 8 is about a famous problem of concurrency. Read this introduction first.

\section{The producer-consumer problem}

The \textbf{producer-consumer problem}, also called the \textbf{bounded buffer problem}, is one of the classic
problems of concurrent programming. Edsger Dijkstra described it in 1965.

Some threads \textbf{produce} things (the \textbf{producers}), other threads \textbf{use} them (the
\textbf{consumers}). They work at different speeds, so between them there is a \textbf{buffer}: a queue that holds
at most $N$ things.

\begin{center}\bufferfig\end{center}

There are two kinds of waiting:
\begin{itemize}
  \item \textbf{The buffer is full:} the producer waits until the consumer takes something.
  \item \textbf{The buffer is empty:} the consumer waits until the producer adds something.
\end{itemize}

\textbf{With mutexes only, the program is correct}: the mutex protects the queue, so nothing is lost or broken.
But a mutex cannot put a thread to sleep until something changes. A thread that waits for a full or an empty buffer
has to unlock, lock and look again in a loop (busy waiting), and it uses the processor for nothing.
\textbf{Condition variables do not fix a bug here, they make the waiting cheap}: the result is the same, but the
waiting thread sleeps. You need three things:

\begin{center}
\renewcommand{\arraystretch}{1.2}
\begin{tabularx}{\linewidth}{@{}lX@{}}
\toprule
\textbf{a mutex} & protects the queue, because both threads change it \\
\textbf{\texttt{not\_full}} & a condition variable: the producer waits on it when the buffer is full; the consumer
  wakes it up after \texttt{pop()} \\
\textbf{\texttt{not\_empty}} & a condition variable: the consumer waits on it when the buffer is empty; the producer
  wakes it up after \texttt{push()} \\
\bottomrule
\end{tabularx}
\end{center}

\textbf{Why a buffer at all?} It smooths out the differences in speed. When the producer is slow for a moment,
the consumer takes from the stock in the buffer. When the consumer is slow for a moment, the producer puts things
into the buffer and does not have to wait.

This pattern is everywhere: a video player, a printer queue, a web server that handles requests, and the pipe in
a terminal (\texttt{cat file | grep word}: \texttt{cat} produces lines, \texttt{grep} consumes them).

\begin{keymsg}
Producer-consumer: two kinds of waiting (full, empty) around one buffer, protected by one mutex.
\end{keymsg}

\section{The task}

\begin{story}
You work for a company that writes a driver for video decoding. One thread decodes the frames of a video (the
producer), and another one shows them on the screen at a fixed speed: one frame every 40\,ms, 25 frames per
second (the consumer). Decoding takes an uneven time: a normal frame is fast, a key frame (an I-frame) is much
slower. Between the threads there is a buffer of decoded frames. When the buffer is empty, the picture freezes.
The driver works correctly, but customers complain that their laptops get hot and the battery runs down fast.
Your team must fix it.
\end{story}

\begin{center}
\begin{tikzpicture}
  % the decoder (producer)
  \node[actor=sIndigo, minimum width=2.6cm, minimum height=1.5cm, align=center, font=\small\bfseries] (dec) at (0,0)
    {\Large\faFilm\\[2pt] decoder};
  \node[font=\scriptsize, text=nInk, align=center, anchor=north] at (0,-0.95)
    {the producer\\normal frame: 5\,ms\\key frame (every 25th): 200\,ms};
  % the buffer of decoded frames
  \foreach \k in {0,...,5} \node[cellused, minimum size=0.75cm, font=\scriptsize] at (2.6+\k*0.78,0) {\#\k};
  \foreach \k in {6,...,9} \node[cellfree, minimum size=0.75cm] at (2.6+\k*0.78,0) {};
  \draw[decorate, decoration={brace, amplitude=5pt, mirror}, nInk] (2.2,-0.55) -- (9.98,-0.55)
    node[midway, below=6pt, font=\scriptsize, text=nInk, align=center] {the buffer: at most \texttt{CAPACITY} = 10 decoded frames\\(\texttt{std::queue<Frame>})};
  \draw[data] (dec.east) -- node[above, font=\scriptsize] {push} (2.18,0);
  % the display (consumer)
  \node[actor=sPink, minimum width=2.6cm, minimum height=1.5cm, align=center, font=\small\bfseries] (scr) at (12.0,0)
    {\Large\faTv\\[2pt] display};
  \draw[data] (10.02,0) -- node[above, font=\scriptsize] {pop} (scr.west);
  \node[font=\scriptsize, text=nInk, align=center, anchor=north] at (12.0,-0.95)
    {the consumer\\one frame every 40\,ms\\empty buffer: the picture freezes};
\end{tikzpicture}
\end{center}

\begin{note}{Ready code (in \texttt{video.h}, we do not look inside)}
\textbf{Functions}\\[2pt]
\begin{tabularx}{\linewidth}{@{}p{3.6cm}X@{}}
\texttt{decode\_frame(i)} & decode frame \texttt{i}: a normal frame takes 5\,ms, every 25th frame is a key frame and takes 200\,ms \\
\texttt{show(f)} & show the frame for 40\,ms; if it comes more than 10\,ms late, print that the picture froze \\
\texttt{report()} & print how many times the picture froze \\
\end{tabularx}
\par\smallskip\textbf{Constants}\\[2pt]
\begin{tabularx}{\linewidth}{@{}p{3.6cm}X@{}}
\texttt{FRAMES} & 100 frames: the video lasts 4\,s \\
\end{tabularx}
\end{note}

The template is in \ghfolder{Task8}. This is \texttt{video.cpp}:

\codefile[firstline=7]{video.cpp}

\runline

\section{Your tasks}

\begin{questions}
  \setcounter{questionsi}{-1}
  \item Who is the producer here, who is the consumer, and what is the buffer? When does each of them wait?
  \item Run it with \texttt{time ./build/video}. How much processor time (\texttt{user}) do the threads use?
        (The program does not end: stop it with Ctrl+C about 2\,s after \texttt{decoder: all frames decoded}.
        Question 4 fixes this.)
  \item Replace busy waiting with two condition variables, \texttt{not\_empty} and \texttt{not\_full}.
        Who waits on which one, and who wakes up whom? Run it with \texttt{time} again.
  \item A colleague put \texttt{std::this\_thread::sleep\_for(std::chrono::milliseconds(10))} into the waiting
        loops instead. The laptop stays cool. Is this as good a solution as a condition variable?
  \item The video is over, but the program does not end. Why? Fix it, so that the display shows all frames that
        are still in the buffer and then ends. (Hint: the decoder must tell the display that nothing more comes.)
  \item Set \texttt{CAPACITY} to 1, then 10, then 100. How many times does the picture freeze?
        Why does the decoder work ahead? What does a larger buffer cost?
\end{questions}
